
\documentclass[dvipdfmx]{jsarticle}

\usepackage[dvipdfmx]{hyperref}
% *** usepackage (アルファベット順) ***
\usepackage{
    amsfonts,amssymb,amsthm,
    braket,
    enumitem,
    physics,
    tikz
}
\newenvironment{probblock}[2]{
  \vspace{0.35\baselineskip} %上余白調節
    %
    % *** 問題番号・点数 ***
    \hspace{-1.5zw} \probbox{#1} \hfill 
      \raise 0.3\baselineskip \hbox{(#2)} \par
    %
  \vspace{-0.15\baselineskip} %下余白調節
}{}

\newcommand{\probbox}[1]{
  \begin{tikzpicture}
    \def\HALF{0.42}
    \coordinate (A) at (-\HALF,-\HALF);
    \coordinate (B) at (\HALF,-\HALF);
    \coordinate (C) at (\HALF,\HALF);
    \coordinate (D) at (-\HALF,\HALF);
    \draw (A)--(B)--(C)--(D)--cycle;
    \node at (0,0) {\huge #1};
  \end{tikzpicture}
}

% [*** 使い方 ***]
% \begin{probblock}{問題番号}{点数}
%   問題文
% \end{probblock}
%
% [*** 例 ***]
% \begin{probblock}{1}{334点}
%   $\tan 1^{\circ}$ は有理数か。
% \end{probblock}
%
% [*** 蛇足 ***]
% 個人的に見栄えがよさげな余白に調節しています。
% もし気に入らない感じだったら変更してください。
% (by ぷりん)


% *** 余白調整 ***
\usepackage[%
top    = 30truemm,%
bottom = 20truemm,%
left   = 25truemm,%
right  = 25truemm]{geometry}

\newtheorem{theo}{定理}[section]
\newcommand{\theocount}[2]{\setcounter{subsection}{#1}\setcounter{theo}{#2}}
\newcommand{\propcount}[2]{\setcounter{subsection}{#1}\setcounter{prop}{#2}}
\newtheorem{prop}[theo]{命題}
\newtheorem{hint}{Hint}[section]
\newtheorem{prob}{問題}[section]
\newtheorem{tips}[theo]{Tips}
\newtheorem{cor}{系}[theo]
\newtheorem{defn}[theo]{定義}
\newtheorem{eg}[theo]{例}
\newtheorem{egprob}[theo]{例題}
\newtheorem{egans}{解答}[theo]
\newtheorem{rem}[theo]{注意}
\newtheorem{rev}{復習}
\newtheorem{assume}{仮定}[theo]
\newtheorem{ans}{解答}[theo]
\newtheorem{que}{疑問}
\newtheorem*{recall}{Recall}
\newtheorem*{theo*}{定理}
\newtheorem*{prop*}{命題}
\newtheorem*{claim*}{Claim}
\newtheorem*{lem*}{補題}
\newtheorem*{defn*}{定義}
\newtheorem*{cor*}{系}
\newtheorem{probans}{解答}[section] %%演習問題用
\newtheorem*{prob*}{問題}
 %newtheorem

%Enumerate 環境
\def\theenumi{\arabic{enumi}}
\def\labelenumi{(\theenumi)}

% (Re)New commands

\definecolor{shadecolor}{gray}{0.80}


\newcommand{\bbox}{\quad\rule{2.2mm}{3mm}}
\newcommand{\aran}{\leftarrow}
\newcommand{\step}[2]{\textbf{Step #1.} (#2)\quad} 
\newcommand{\dothom}[1]{ \underset{#1}{\overset{\cdot}{\simeq}}  }
\newcommand{\maru}{\textcircled}  %%%まるで囲む数字
\newcommand{\prf}{\noindent{\large (\textbf{Proof.})\quad}} 
\newcommand{\ten}{$\bullet\quad$} %黒丸
\newcommand{\vin}{\rotatebox{90}{$\in$}}
\newcommand{\thm}{\begin{itembox}[l]}
\newcommand{\enthm}{\end{itembox}\\}
\newcommand{\disp}{\displaystyle}
\newcommand{\dou}{\Leftrightarrow}
\newcommand{\ben}{\begin{enumerate}}
\newcommand{\een}{\end{enumerate}}
\newcommand{\beqn}{\begin{eqnarray*}}
\newcommand{\eeqn}{\end{eqnarray*}}
\newcommand{\bcas}{\begin{cases}}
\newcommand{\ecas}{\end{cases}}
\newcommand{\spec}[1]{\mr{Spec}\parena{#1}}
\newcommand{\nara}{\Rightarrow}
%%%フォント
\newcommand{\mb}{\mathbb} %%白抜き
\newcommand{\mr}{\mathrm} %%%ローマン
\newcommand{\maf}{\mathfrak} %%%フラクトゥール(素イデアルなど)
\newcommand{\hana}{\mathscr} %%%花文字(圏のやつ)
\newcommand{\mac}{\mathcal} %%%関手のやつ
\newcommand{\besha}{\begin{shadebox}}
\newcommand{\ensha}{\end{shadebox}}

\newcommand{\chat}{\widehat{\mathbb{C}}} %リーマン球面
%\newcommand{\khat}{\widehat{K}}
\newcommand{\nea}{\nearrow}
\newcommand{\bolm}[1]{{\mbox{\boldmath $#1$}}}
\newcommand{\exer}{{\bf 演習}\arabic{chapter}.\arabic{section}.}
\newcommand{\sea}{\searrow}
\newcommand{\tf}{\textbf}

\newcommand{\bracket}[1]{%
\[
 \left(
 \begin{tabular}{p{0.9\hsize}}
  #1
 \end{tabular}
 \right)
\]}

\renewcommand{\leq}{\leqq}
\renewcommand{\geq}{\geqq}
\renewcommand\thefootnote{*\arabic{footnote}}
\newcommand{\kuro}[1]{\mr{\textbf{#1}}}  %%複体の文字とか

%%%代数学
\newcommand{\tens}{\mathbin{\otimes}} %%%テンソル
\newcommand{\ilim}[1][]{\mathop{\varinjlim}\limits_{#1}}%%逆極限
\newcommand{\plim}[1][]{\mathop{\varprojlim}\limits_{#1}}%順極限
\newcommand{\im}[1]{\mr{Im}\,#1}
\newcommand{\coker}[1]{\mr{Coker}\,#1}
\newcommand{\rank}[1]{\mr{rank}\, #1}
\newcommand{\tor}[2]{\mr{Tor}_{#1}^{#2}}  %%Tor{i番目}{環}
\newcommand{\ext}[2]{\mr{Ext}^{#1}_{#2} }
\renewcommand{\hom}[1]{\mr{Hom}\parena{#1}} %Hom関手
\newcommand{\norgr}{\rotatebox{90}{$\triangle$}}  %%正規部分群
\newcommand{\bep}{\begin{pmatrix}}
\newcommand{\enp}{\end{pmatrix}} %行列
\newcommand{\tr}[1]{\mr{Tr}\parena{#1}} %Trace
\newcommand{\gal}[1]{\mr{Gal}\parena{#1}} %Galois群

%%%幾何学
\newcommand{\supp}[1]{\mr{Supp}(#1)}
\newcommand{\opensub}{\overset{\mr{open}}{\subset}}


%%%解析学
\newcommand{\grad}{\mr{grad}} %%%grad
\newcommand{\inpr}[2]{\lan #1, #2 \ran} %%%内積
\newcommand{\del}{\partial} %%%偏微分, 境界
\newcommand{\abs}[1]{ \left| #1 \right| }
\newcommand{\norm}[1]{\left\| #1 \right\| }
\renewcommand{\span}[1]{\mr{Span}(#1)}
\newcommand{\fou}[1]{\mac{F}\parenc{#1}}
\newcommand{\llfou}[1]{\mac{F}_{L^2}\parenc{#1}}
\newcommand{\uof}[1]{\mac{F}^{-1}\parenc{#1}}
\newcommand{\lluof}[1]{\mac{F}^{-1}_{L^2}\parenc{#1}}
\newcommand{\conv}[1]{\mr{Conv}\, (#1)}
\newcommand{\res}[2]{\underset{\mr{Res}}{#1}#2}
\newcommand{\slim}{\mr{s-}\disp\lim}
\newcommand{\wlim}{\mr{w-}\disp\lim}


%%%その他数式用
\newcommand{\witi}{\widetilde} %%%ワイドチルダ
\newcommand{\wiha}{\widehat} %%%ワイドはっと(フーリエ係数)
\newcommand{\symdiff}{\bigtriangleup}  %%%対称差
\newcommand{\ovl}{\overline}
\newcommand{\lan}{\langle}
\newcommand{\ran}{\rangle}

\newcommand*{\ve}{\overrightarrow} %%%%%ベクトル
\newcommand{\parena}[1]{\left( #1 \right)}
\newcommand{\parenb}[1]{\left\{ #1 \right\}}
\newcommand{\parenc}[1]{\left[ #1 \right]}
\newcommand{\bea}[2]{\begin{array}{#1} #2 \end{array}}


\newcommand{\reg}[1]{\mr{Reg}\, #1}  %%非特異点集合
\renewcommand{\div}[1]{\mr{div}(#1)}
\newcommand{\sgn}[1]{\mr{sgn}#1}


%%%左赤線環境
\makeatletter

\renewenvironment{leftbar}{%
%  \def\FrameCommand{\vrule width 3pt \hspace{10pt}}%  デフォルトの線の太さは3pt
  \def\FrameCommand{\vrule width 1pt \hspace{10pt}}% 
  \MakeFramed {\advance\hsize-\width \FrameRestore}}%
 {\endMakeFramed}
\newcommand{\lefba}[1]{\begin{leftbar} #1 \end{leftbar}}

\makeatletter

\def\mapstofill@{%
   \arrowfill@{\mapstochar\relbar}\relbar\rightarrow}
\newcommand*\xmapsto[2][]{%
   \ext@arrow 0395\mapstofill@{#1}{#2}}
\newcommand{\xequal}[2][]{\ext@arrow 0055{\equalfill@}{#1}{#2}}
\def\equalfill@{\arrowfill@\Relbar\Relbar\Relbar}
\makeatother


% 問題文用の環境 'probblock' を作りました。
% 定義は __def_probblock_env.tex に書いてあります。
% (by ぷりん)
\include{__def_probblock_env}

% *** ページ番号を消去 ***
\pagestyle{empty}

% *** フォントサイズ・行送りの指定 ***
\newcommand{\applyfontsetting}{
  \fontsize{11pt}{17pt}\selectfont
  \renewcommand{\baselinestretch}{1.05}
}

% @@@@@@@@@@@@@@@@@@@@@@@@@@@@@@@@@@@@@@@@@@@@@@@@@@@@@@
% @@@@@@@@@@@@@@@@@@@@@@@@@@@@@@@@@@@@@@@@@@@@@@@@@@@@@@
\begin{document}
\applyfontsetting

\section{第1セット(数学I)}

% *************
% ***** 1 *****
% *************
\begin{probblock}{1}{H31東北大院試(共同)}
3次実正方行列
\[A = \bep  0 & -1 & 1 \\ -1 & 0 & -1 \\ 1 & -1 & 0 \enp\]
について, 以下の問いに答えよ。
\ben
\item $A$の固有値をすべて求めよ。また, 各固有値に対応する固有空間の基底を一組求めよ。
\item 極限$\disp\lim_{n\to \infty} \dfrac{1}{2^{n}}A^{n}\bep 5 \\ 0 \\ 1 \enp$ を求めよ。
\een
\end{probblock}

\vfill

% *************
% ***** 2 *****
% *************
\begin{probblock}{2}{微分積分(吉田伸夫)問12.2.3}
$f:[0,\infty)\to [0,\infty)$は狭義単調減少な連続関数とする。二つの広義積分
\[I=\int_{0}^{\infty} f(x)\, dx,\qquad J=\int_{0}^{\infty} f(x)|\sin{x}|\, dx\]
に対し, $I<\infty$であることと$J<\infty$であることは同値であることを証明せよ。
\end{probblock}

\vfill

% *************
% ***** 3 *****
% *************
\begin{probblock}{3}{\href{http://zen.shinshu-u.ac.jp/modules/0023000000/}{環論演習第11回}}
$\mb{Z}[x]$の部分環$R$を$R = \mb{Z}[2x,2x^2, 2x^3,\cdots⋅]$とする。即ち, 多項式$f(x) = a_nx^n +\cdots + a_1x + a_0 \in \mb{Z}[x]$について，$f(x)\in R \overset{\mr{def}}{\iff}  \text{$a_1,\cdots ,a_n$は 2の倍数 ($a_0$は任意の整数)}$である。
$R$の部分集合$\maf{a}$を
\[\maf{a} = \{ a_nx^n + \cdots⋅ + a_1x \mid n \geq 1, a_1\cdots ,a_n\text{は2の倍数} \} \]
とおく。このとき
\ben
\item $\maf{a}$ は $R$のイデアルであることを示せ。
\item $R$はNoether環か?
\een
\end{probblock}
\vfill

% *************
% ***** 4 *****
% *************
\begin{probblock}{4}{(2020幾何学I問題65 改)}
$S^2 = \{ (x,y,z)\mid x^2+y^2+z^2=1 \}$上の関数$f:S^2\to \mb{R}$を$f(x,y,z)=xy$で定める。$f(x,y,z)$の臨界点をすべて求めよ。
\end{probblock}

\vfill

% *************
% ***** 5 *****
% *************
\begin{probblock}{5}{複素関数概論, 定理4.7}
$f$を整関数とする。整数$p\geq 0$, 定数$M$および$\lim_{n\to \infty}r_n = \infty$をみたす正数の列$\{ r_n \}_{n=1}^{\infty}$が存在して, 
\[\max_{|z| = r_n} |f(z)| \leq Mr_n^{p}\quad (n=1,2,\ldots)\]
が成り立てば, $f$は高々$p$次の多項式であることを証明せよ。
\end{probblock}
\vfill


\clearpage

%全滅系問題
%H21 なし
%H22 4番(逆関数定理)
%H23 1,4番
%H24 4(写像度2) , 3(F_5イデアルx^4+2x^2)
%H25 3(X^pが単位行列), 4(部分多様体)
%H26 なし
%H27 2(一様収束)  4(臨界点)
%H28 3(単院試論)
%H29 4(臨界点),5(最大値原理)

\section{第2セット(ほぼ全滅系問題)}
留学生入試の完答者1人以下で, 教育的な問題のセット。少し誘導を追加したので復習しよう。



\begin{probblock}{1}{H25}
$\mathbb{F}_{p}$を成分に持つ2×2正方行列$X$で, $X^p = I$を満たすものがいくつあるかについて調べたい。
\ben
\item $a\in \ovl{\mathbb{F}_{p}} = \bigcup_{k\geq 1} \mb{F}_{p^k}$が$a^{p} = 1$を満たすならば$a = 1$であることを示せ。
\item $X$の固有値を$\lambda_1,\lambda_2$とする。$\lambda_i$は$\mb{F}_{p}$に属すとは限らないことに注意せよ。 $\lambda_i \in \ovl{\mathbb{F}_{p}} = \bigcup_{k\geq 1} \mathbb{F}_{p^k}$である。$X^p = I \iff \lambda_i =1 (i=1,2)$を示せ。
\item 結局, 固有値が1のみであるような$X$の個数を求めればよい。$X$は何個あるか。
\een
\end{probblock}

\begin{probblock}{2}{H27}
$\{ f_n \}$を開区間$(0,1)$上の微分可能な関数の列で, 任意の$x\in (0,1)$と任意の正の整数$n$に対して$|f_n'(x)|\leq 1$が満たされているとする。$\{ f_n \}$が$f$に$(0,1)$で各点収束するとする。
\ben 
\item $|f(x) - f(y)| \leq |x-y|$を示せ。
\item 任意の$\epsilon > 0$に対し, $f_n$は閉区間$[\epsilon,1-\epsilon]$上一様収束することを示せ。
\item $f_n$は$[0,1]$上一様収束することを示せ。 
\een
\end{probblock}

\begin{probblock}{3}{(1):中国剰余定理 (2)H24 (3)H27 (4)H26 }
\ben 
\item $A$が単位元を持つ可換環で, $I,J$が$I+J = A$を満たすイデアルとする。このとき, $IJ = I\cap J$であることを示し, 次は環同型であることを証明せよ。
\[
A/IJ \cong A/I \times A/J,\quad a\bmod{IJ} \leftrightarrow (a\bmod{I}, b\bmod{J})
\]
\item 剰余環$\mb{F}_{5}[x]/(x^4+2x^2)\mb{F}_{5}[x]$の可逆元の個数を求めよ。
\item $\mathbb{R}[x]/(x^3-8)$の元$a$で$a^4 = 1$を満たすものの個数を求めよ。
\item 環準同型$\phi:\mb{C}[x]/(x^2(x+1))\to \mb{C}$で$a\in \mb{C}$ならば$\phi(a)=a$であるようなものをすべて求めよ。
\een
\end{probblock}

\begin{probblock}{4}{H29}
$n$を正の整数, $M$を境界のないコンパクト$n$次元多様体とする。$f:M\to \mb{R}$を非退化な臨界点しかもたない$C^{\infty}$級関数とし, $f$の臨界点の集合を$A$とする。
\ben 
\item 臨界点$p\in A$が非退化であるとはどういうことを言うのか説明せよ。
\item 座標近傍$(U_i;x_1,\cdots, x_n)$において, $g(p) = (\dfrac{\del f}{\del x_1}(p),\cdots, \dfrac{\del f}{\del x_n}(p))$ として$C^{\infty}$級関数$g:U_i\to \mb{R}^{n}$を定義する。$p\in A$ならば$(dg)_{p}:T_{p}M \to \mb{R}^{n}$は線形同型であることを示せ。
\item $U_i\cap A = g^{-1}(\ve{0})$を示し, これが離散集合であることを示せ。そして$A$が有限集合であることを証明せよ。
\een
\end{probblock}
\newpage 
\begin{probblock}{5}{H29}
$D = \{ |z| < 1 \}$. $f,g$はともに$D$上正則かつ$\ovl{D}$上で連続とし, 
\[|f(z)| \leq |g(z)|\, (\mr{on} D),\quad |f(z)| = |g(z)|\, (\text{on} \del D),\quad f(z)\neq 0\, (\text{on} \ovl{D}-\{ 0 \})\]
\ben 
\item $D-\{ 0\}$上の正則函数$h(z) = \frac{f(z)}{z^mg(z)}$は$\ovl{D}$上連続で$D$上正則な関数に拡張できることを示せ。
\item $h^{-1}$は$D$上正則であり,  絶対値の最大値が$1$であることを示せ。そして$|h|\equiv 1$が$\ovl{D}$上で成立することを示せ。
\item $h$は絶対値1の定数であることを示せ。よって, $f(z) = \alpha z^m g(z)$ ($|\alpha| = 1$) と書ける。
\een
\end{probblock}

\begin{probblock}{6}{おまけ(剰余環の練習)}
次の同型を示せ。
\ben
\item $\mb{R}[x]/(x^3+1) \cong \mb{R}\times \mb{C}$
\item $\mathbb{C}[x,y]/(x^2+y^2)\cong \mb{C}\times \mb{C}$
\item $\mb{F}_{19}[x]/(x^2-2) \cong \mb{F}_{361}$
\item $\mb{Z}[x]/(17,x^2-2) \cong \mb{F}_{17}\times \mb{F}_{17}$
\item $\mb{R}[x]/(x^4 + 1) \cong \mb{C}\times \mb{C}$
\item $\mb{Z}[\sqrt{-1}]/(p)\cong \bcas
\mb{Z}/4\mb{Z}\quad (p=2) \\ 
\mb{F}_{p}\times \mb{F}_{p} \quad (p\equiv 1\pmod{4})\\
\mb{F}_{p^2} \quad (p\equiv 3\pmod{4})
\ecas$; $p$は素数
\een
\end{probblock}

\clearpage
\section{第3回(代数題)}


\clearpage
\section{第4回(部誌用)}


\end{document}












\end{document}